\documentclass[11pt]{article}
\usepackage{coursenotes}
\begin{document}
\handoutheader{H4}{Instrumental Variables}{Outside shifters of treatment: instruments that nobody randomised}{Meeting 7 (Observational data)}

\begin{decision}
A firm suspects an unmeasured confounder that no adjustment can fix (Meeting 7), and it cannot run an experiment. Something
outside its control may still shift the treatment: a rule, a manager's choice, a system outage, the weather. Meeting~2 showed
what the ratio of two differences identifies when the shifter is a randomised offer. This handout is about shifters that
nobody randomised: what must then be argued, what goes wrong when the argument fails, and how to tell a strong instrument
from a weak one.
\end{decision}

\begin{goals}
  \item say what changes when the instrument is not randomised, and argue independence and exclusion from the institution;
  \item describe the compliers, the units whose effect the ratio measures;
  \item quantify what failures of exclusion and monotonicity do to the ratio;
  \item interpret two-stage least squares with covariates, and guard against weak instruments.
\end{goals}

%=====================================================================================
\sectionone

\subsection{From a randomised offer to an outside shifter}

The notation is Meeting~2's (Section~1.9): a binary instrument $Z$, a treatment $D$ with potential treatments $D(0), D(1)$,
and potential outcomes $Y(d, z)$. Meeting~2 took $Z$ to be a randomised offer, so independence came from the draw. Here
$Z$ is any variable that shifts $D$, and independence becomes an assumption.

\begin{assum}{Independence}{ind}
$\{Y(d,z), D(z) : d, z \in \{0,1\}\} \indep Z$, possibly only within levels of covariates $X$.
\end{assum}
\begin{assum}{Relevance}{rel}
$\E[D(1) - D(0)] \ne 0$.
\end{assum}

Exclusion and monotonicity are as in Meeting~2 (Assumptions~1.11 and~1.12), with $Z$ in place of the draw. Under all four,
Meeting~2's results carry over unchanged: the ITT and the first stage are identified (Result~1.13), the type shares are
identified (Result~1.14), and the ratio $\text{ITT}/\pi$ is the LATE, the average effect on compliers (Result~1.15). What
changes is the status of each assumption. Relevance is checkable. Independence, which a draw delivers by design, is now a
claim about the world, argued like Meeting~7's exchangeability. Exclusion and monotonicity were already claims; with an
outside shifter they are usually harder to defend, because nobody designed the instrument to satisfy them.

\subsection{Arguing for an instrument}

No test confirms an instrument; each of the following can only count against one.
\begin{itemize}
  \item \emph{Where does $Z$ come from?} Name the process that set it and everything that process responded to. Those
  variables belong in $X$, and independence is claimed only within their levels.
  \item \emph{Balance.} Pre-treatment covariates should not differ by $Z$ within levels of $X$. Imbalance is evidence against
  independence; balance is consistent with it, not proof.
  \item \emph{Zero-first-stage groups.} Units the instrument cannot move (ineligible customers, a system that cannot display
  the offer) should show no outcome difference by $Z$. A difference there is a direct effect: evidence against exclusion,
  and an estimate of its size (Result~\ref{prop:violation}).
  \item \emph{Placebo outcomes.} Outcomes the treatment cannot affect, such as anything measured before $Z$, should not move with $Z$.
\end{itemize}

\subsection{Who the compliers are}

\begin{prop}{Complier outcome means}{complier}
Under independence, relevance, exclusion and monotonicity,
$\E[Y(1) \mid C] = \{\E[YD \mid Z=1] - \E[YD \mid Z=0]\}/\pi$ and $\E[Y(0) \mid C] = \{\E[Y(1-D) \mid Z=0] - \E[Y(1-D) \mid Z=1]\}/\pi$.
\end{prop}
\begin{proof}
$\E[YD \mid Z=1] = \pi_A\E[Y(1) \mid A] + \pi_C\E[Y(1) \mid C]$ and $\E[YD \mid Z=0] = \pi_A\E[Y(1) \mid A]$; subtract and divide. The second uses
never-takers in the same way.
\end{proof}

Replacing $Y$ with a covariate describes the compliers, the first step in arguing whether their effect transports. With an
outside shifter the compliers are whoever the shifter happened to move, a group nobody chose.

\subsection{When exclusion or monotonicity fails}

\begin{prop}{Exclusion violation}{violation}
Keep independence, relevance and monotonicity, but let $Y(d,1) = Y(d,0) + \delta$ for every unit. Then
$\text{ITT}/\pi = \text{LATE} + \delta/\pi$.
\end{prop}
\begin{proof}
$Y(D(1),1) - Y(D(0),0) = \{Y(D(1),0) - Y(D(0),0)\} + \delta$; the first part averages to $\pi\,\text{LATE}$.
\end{proof}

The direct effect is divided by the first stage: harmless with a strong instrument, fatal with a weak one. Report the LATE as a
function of $\delta$ and the $\delta$ at which the decision flips.

\begin{prop}{Defiers}{defiers}
Keep independence, exclusion and relevance but allow defiers. Then
$\text{ITT}/\pi = (\pi_C\,\text{LATE}_C - \pi_D\,\text{LATE}_D)/(\pi_C - \pi_D)$.
\end{prop}
\begin{proof}
The proof of the LATE theorem (Meeting~2, Result~1.15) now leaves a defier term $-\pi_D\,\text{LATE}_D$ in the ITT, and the first
stage is $\pi_C - \pi_D$.
\end{proof}

With defiers one weight is negative, and the ratio need not lie between the two groups' effects; it can have the wrong sign.

\subsection{Covariates and two-stage least squares}

If the instrument is random only given covariates $X$ (assignment probabilities differ by stratum), two-stage least squares with
stratum indicators computes $\hat\theta = \sum_i\tilde Z_iY_i/\sum_i\tilde Z_iD_i$ with $\tilde Z_i = Z_i - \hat\E[Z \mid X_i]$.

\begin{prop}{What 2SLS with strata estimates}{tsls}
Under independence and relevance within each level of $X$, with exclusion and monotonicity, and with discrete $X$,
\[
  \hat\theta \to \frac{\E[\Var(Z \mid X)\,\pi(X)\,\text{LATE}(X)]}{\E[\Var(Z \mid X)\,\pi(X)]}.
\]
\end{prop}
\begin{proof}
For binary $Z$, $\E[\tilde ZY] = \E[\Var(Z \mid X)\,\text{ITT}(X)]$ and $\E[\tilde ZD] = \E[\Var(Z \mid X)\,\pi(X)]$; within each stratum
$\text{ITT}(X) = \pi(X)\,\text{LATE}(X)$.
\end{proof}

Controls that determine the instrument's probability are required; controls that only predict $Y$ reduce variance; anything
measured after assignment is never a control. With heterogeneous effects, 2SLS weights strata towards balanced instruments and
strong first stages. Running the two stages by hand reproduces the estimate but not its standard error; use an IV routine. With a
continuous treatment (a log price) the same formula gives an elasticity, a weighted average of slopes over units whose treatment
the instrument moved.

\subsection{Inference and weak instruments}

\begin{prop}{Delta-method standard error}{delta}
With $\hat\theta = \widehat{\text{ITT}}/\hat\pi$, approximately
$\Var(\hat\theta) \approx \{\Var(\widehat{\text{ITT}}) - 2\theta\Cov(\widehat{\text{ITT}}, \hat\pi) + \theta^2\Var(\hat\pi)\}/\pi^2$.
\end{prop}
\begin{proof}
Linearise $a/b$ at $(\text{ITT}, \pi)$: $a/b - \theta \approx \{(a - \text{ITT}) - \theta(b - \pi)\}/\pi$.
\end{proof}

With a precise first stage, $\SE(\hat\theta) \approx \SE(\widehat{\text{ITT}})/\pi$. The approximation fails when $\pi$ is small relative to its
SE, measured by the first-stage $F = (\hat\pi/\SE(\hat\pi))^2$. A weak instrument's estimate is heavy-tailed, centred between the truth
and the naive regression, and its nominal 95\% interval undercovers. $F > 10$ (Staiger and Stock) is a floor, not a certificate: Lee,
McCrary, Moreira and Porter (2022) show a conventional 5\% $t$-test has correct size only for $F$ above about $104.7$.

\begin{prop}{The Anderson--Rubin test}{ar}
If $Y = \theta D + u$ with $Z \indep u$, the regression of $Y - \theta_0D$ on $Z$ has slope zero exactly when $\theta_0 = \theta$ (given $\pi \ne 0$),
and the usual test of that slope has correct size whatever the first stage's strength.
\end{prop}
\begin{proof}
$Y - \theta_0D = u + (\theta - \theta_0)D$, whose covariance with $Z$ is $(\theta - \theta_0)\Cov(Z, D)$; under the null the dependent variable is
$u$, independent of $Z$.
\end{proof}

The set of $\theta_0$ not rejected is a confidence set valid under weak instruments; it may be very wide or unbounded, and that
width is an accurate summary of what the data can say. Decide the reporting rule before computing the estimate: report the first stage and $F$; below 10, report the ITT and first stage
rather than the ratio; between 10 and about 100, report the Anderson--Rubin set.

\subsection{Locality}

The LATE averages over compliers, defined by the instrument; a different instrument has different compliers and, with heterogeneous
effects, a different LATE. Ask who the compliers are (Result~\ref{prop:complier}), how large a change in treatment the instrument
induced, and whether the decision is about the compliers at all.

\begin{pitfall}[IV errors]
Treating an outside shifter as random without naming what set it; leaving out the covariates that determined the instrument;
dividing a small ITT by a weak first stage and reporting the ratio; controlling for anything measured after the instrument; reading a
LATE as the ATE; ignoring a direct effect of the instrument.
\end{pitfall}

%=====================================================================================
\sectiontwo

\paragraph{Setting.} QuickBite wants the effect of its Plus membership ($D$) on orders ($Y$, orders in the following month).
Each order earns ¥10 of margin; each Plus member costs ¥15 a month in waived delivery fees. Members choose to join, and heavy
users join most, so comparing members with non-members is hopeless. In March, regional managers in 6 of the 12 cities chose
to run an in-app campaign promoting Plus ($Z = 1$ for a customer in a campaign city). Nobody randomised the campaign: managers
in big cities were keener (4 of 6 big cities ran it, 2 of 6 small ones), and big-city customers order more.

\begin{center}
\small
\begin{tabular}{lrrrrrr}
\toprule
& \multicolumn{3}{c}{Campaign city ($Z = 1$)} & \multicolumn{3}{c}{No campaign ($Z = 0$)} \\
\cmidrule(lr){2-4}\cmidrule(lr){5-7}
City size & Customers & Plus & Orders & Customers & Plus & Orders \\
\midrule
Big & 8{,}000 & 0.30 & 5.30 & 4{,}000 & 0.05 & 5.00 \\
Small & 2{,}000 & 0.30 & 3.30 & 6{,}000 & 0.05 & 3.00 \\
\bottomrule
\end{tabular}
\end{center}

\paragraph{Step 1: the pooled ratio.} Pooled, orders are $4.90$ in campaign cities and $3.80$ elsewhere, an ITT of $1.10$; Plus
membership is $0.30$ against $0.05$, a first stage of $0.25$. The ratio $1.10/0.25 = 4.4$ extra orders per complier is not a
LATE: campaign cities are mostly big cities, so the ITT mixes the campaign's effect with the gap between big and small cities.
Independence fails unconditionally.

\paragraph{Step 2: within city size.} City size determined the campaign's probability, so it belongs in $X$. Within each size
the ITT is $0.30$ and the first stage $0.25$, so the ratio is $1.2$ in both, and two-stage least squares with city-size
indicators (Result~\ref{prop:tsls}) returns $1.2$, with an SE of about $0.22$ for an outcome SD of $3.5$. The first stage is
strong: $F$ is about 1{,}600 in big cities and 550 in small ones. Within city size, independence is still a claim, that managers
did not choose the campaign because they expected their city's orders to rise. Balance on pre-March orders within each size
is the first check.

\paragraph{Step 3: who the compliers are.} In big cities, $\E[YD \mid Z=1] = 2.15$, $\E[YD \mid Z=0] = 0.40$,
$\E[Y(1-D) \mid Z=0] = 4.60$ and $\E[Y(1-D) \mid Z=1] = 3.15$. By Result~\ref{prop:complier}, compliers would order
$(4.60 - 3.15)/0.25 = 5.8$ times a month without Plus and $(2.15 - 0.40)/0.25 = 7.0$ with it. Always-takers, who joined
without a campaign, order $0.40/0.05 = 8.0$; never-takers $3.15/0.70 = 4.5$. The campaign moved middling customers, not the
heaviest ones.

\paragraph{Step 4: exclusion.} The campaign cities also put up billboards. Corporate accounts cannot join Plus, so for them the
first stage is zero, and within each city size they ordered $0.05$ more in campaign cities: a direct effect $\delta = 0.05$.
By Result~\ref{prop:violation} the ratio overstates the LATE by $0.05/0.25 = 0.2$; the corrected complier effect is
$(0.30 - 0.05)/0.25 = 1.0$.

\paragraph{Step 5: the decisions.}
\begin{itemize}
  \item \emph{Run the campaign, billboards and all, in every city?} The within-size ITT is the policy effect, billboards
  included: per customer, $10 \times 0.30 = $ ¥3.00 of margin against $0.25 \times 15 = $ ¥3.75 of waived fees for the extra
  members. Net $-$¥0.75: no.
  \item \emph{Is Plus worth it for the customers the campaign recruits?} $10 \times 1.0 = $ ¥10 a month against ¥15 (and
  $10 \times 1.2 = $ ¥12 before the exclusion correction): no, for compliers. The data say nothing about always-takers, the
  heaviest users, for whom Plus may well pay.
\end{itemize}

\begin{exercise}[computation]
Suppose instead the campaign had lifted membership only from $0.05$ to $0.07$ in each city size ($\pi = 0.02$, $\SE(\hat\pi) \approx 0.0034$),
with the same within-size ITT of $0.30$ (SE $0.05$). Compute the first-stage $F$, the ratio and its approximate SE, and say what you would
report.
\end{exercise}
\answer{$F = (0.02/0.0034)^2 \approx 35$. Ratio $0.30/0.02 = 15$ orders per complier, SE roughly $0.05/0.02 = 2.5$ before the first-stage term,
and more after it. An effect of 15 orders a month per complier is implausible: any direct effect of the campaign of $\delta$ is multiplied by
$1/0.02 = 50$. With $F$ between 10 and 100, report the ITT and first stage, plus an Anderson--Rubin set; do not headline the ratio.}

\begin{exercise}[judgement]
An analyst proposes using rainfall as an instrument for QuickBite's delivery fee (rain reduces courier supply, which raises the
surge fee) to estimate the fee elasticity of orders. Assess each of the four assumptions.
\end{exercise}
\answer{Relevance: checkable, likely strong. Independence: rain is not chosen by anyone, but it is seasonal and regional; compare within
city and season. Monotonicity: rain raises fees everywhere, plausibly. Exclusion: almost certainly fails, because rain raises demand for
delivery directly (people stay in), so the ratio mixes the fee effect with the rain effect on demand, biasing the elasticity upward
(towards zero or positive). A better instrument shifts the fee without shifting demand, such as randomised fee tests or courier-pay
changes.}

%=====================================================================================
\sectionthree

\begin{extension}[Continuous treatments, flexible controls and describing compliers]
With a continuous treatment and binary instrument, the Wald ratio is a weighted average of unit slopes over the range the instrument
moved them (Angrist, Graddy and Imbens, 2000). With many controls, residualise $Y$, $D$ and $Z$ on $X$ with cross-fitted learners and compute
$\sum\tilde Z\tilde Y/\sum\tilde Z\tilde D$: the partially linear IV model of double machine learning. Abadie's $\kappa$-weighting identifies any feature of the
compliers' joint distribution.
\end{extension}

\begin{extension}[Sensitivity to exclusion]
Report the LATE as a function of the direct effect, $(\text{ITT} - \delta)/\pi$, over a range of $\delta$ argued from the institution, and the
$\delta$ at which the decision flips. Conley, Hansen and Rossi (2012) formalise this as ``plausibly exogenous'' instruments.
\end{extension}

\subsection*{Annotated reading}
\begin{readinglist}
\reading{Angrist, Imbens and Rubin (1996), ``Identification of causal effects using instrumental variables'', \emph{JASA}}{The
potential-outcomes framework for IV, compliance types and the LATE (also in Meeting~2's list)}{Sections 1--4}{2}
\reading{Angrist and Pischke (2015), \emph{Mastering 'Metrics}, Princeton}{IV at an intuitive level, with classic
examples}{Ch.\ 3}{1}
\reading{Imbens and Angrist (1994), ``Identification and estimation of local average treatment effects'',
\emph{Econometrica}}{The LATE theorem}{All}{3}
\reading{Lee, McCrary, Moreira and Porter (2022), ``Valid t-ratio inference for IV'', \emph{AER}}{Why $F > 10$ is not enough, and
corrected critical values}{Sections I--III}{3}
\reading{Andrews, Stock and Sun (2019), ``Weak instruments in instrumental variables regression'', \emph{Annual Review of
Economics}}{A survey of weak-instrument-robust inference}{Sections 1--4}{3}
\reading{Graddy (2006), ``The Fulton Fish Market'', \emph{Journal of Economic Perspectives}}{Weather as an instrument for price in a
real market}{All}{1}
\end{readinglist}

\end{document}
